\section{USAMO 2019}
        \begin{enumerate}
        	\item (\textit{USAMO 2019}) \textbf{Note:} For any geometry problem whose statement begins with an asterisk $(*)$, the first page of the solution must be a large, in-scale, clearly labeled diagram. Failure to meet this requirement will result in an automatic 1-point deduction.


 Let $\mathbb{N}$ be the set of positive integers. A function $f:\mathbb{N}\to\mathbb{N}$ satisfies the equation 
 \[\underbrace{f(f(\ldots f}_{f(n)\text{ times}}(n)\ldots))=\frac{n^2}{f(f(n))}\]for all positive integers $n$. Given this information, determine all possible values of $f(1000)$.


 

	\item (\textit{USAMO 2019}) \textbf{Note:} For any geometry problem whose statement begins with an asterisk $(*)$, the first page of the solution must be a large, in-scale, clearly labeled diagram. Failure to meet this requirement will result in an automatic 1-point deduction.


 Let $ABCD$ be a cyclic quadrilateral satisfying $AD^2 + BC^2 = AB^2$. The diagonals of $ABCD$ intersect at $E$. Let $P$ be a point on side $\overline{AB}$ satisfying $\angle APD = \angle BPC$. Show that line $PE$ bisects $\overline{CD}$.


 

	\item (\textit{USAMO 2019}) \textbf{Note:} For any geometry problem whose statement begins with an asterisk $(*)$, the first page of the solution must be a large, in-scale, clearly labeled diagram. Failure to meet this requirement will result in an automatic 1-point deduction.


 Let $K$ be the set of all positive integers that do not contain the digit $7$ in their base-$10$ representation. Find all polynomials $f$ with nonnegative integer coefficients such that $f(n)\in K$ whenever $n\in K$.


 



 

	\item (\textit{USAMO 2019}) \textbf{Note:} For any geometry problem whose statement begins with an asterisk $(*)$, the first page of the solution must be a large, in-scale, clearly labeled diagram. Failure to meet this requirement will result in an automatic 1-point deduction.


 Let $n$ be a nonnegative integer. Determine the number of ways that one can choose $(n+1)^2$ sets $S_{i,j}\subseteq\{1,2,\ldots,2n\}$, for integers $i,j$ with $0\leq i,j\leq n$, such that:
for all $0\leq i,j\leq n$, the set $S_{i,j}$ has $i+j$ elements; and
$S_{i,j}\subseteq S_{k,l}$ whenever $0\leq i\leq k\leq n$ and $0\leq j\leq l\leq n$.


 

	\item (\textit{USAMO 2019}) \textbf{Note:} For any geometry problem whose statement begins with an asterisk $(*)$, the first page of the solution must be a large, in-scale, clearly labeled diagram. Failure to meet this requirement will result in an automatic 1-point deduction.


 Two rational numbers $\frac{m}{n}$ and $\frac{n}{m}$ are written on a blackboard, where $m$ and $n$ are relatively prime positive integers. At any point, Evan may pick two of the numbers $x$ and $y$ written on the board and write either their arithmetic mean $\frac{x+y}{2}$ or their harmonic mean $\frac{2xy}{x+y}$ on the board as well. Find all pairs $(m,n)$ such that Evan can write $1$ on the board in finitely many steps.


 

	\item (\textit{USAMO 2019}) \textbf{Note:} For any geometry problem whose statement begins with an asterisk $(*)$, the first page of the solution must be a large, in-scale, clearly labeled diagram. Failure to meet this requirement will result in an automatic 1-point deduction.


 Find all polynomials $P$ with real coefficients such that 
 \[\frac{P(x)}{yz}+\frac{P(y)}{zx}+\frac{P(z)}{xy}=P(x-y)+P(y-z)+P(z-x)\]holds for all nonzero real numbers $x,y,z$ satisfying $2xyz=x+y+z$.


 

\end{enumerate}